\subsection{Sturm-Liouville 边值问题}\label{sec:sturm-liouville-boundary-value-problems}

\PDEIndexMark{idx:sturm-liouville-problem-and-regularity}现在研究由分离变量法产生的一类二阶 ODE 边值问题. \textbf{Sturm-Liouville 问题} (或 S-L 问题) 包含区间 $(a,b)$ 上的常微分方程
\begin{equation}\label{eq:sturm-liouville-inhomogeneous-differential-equation}
-\frac{\mathrm{d}}{\mathrm{d}x}\left(p(x)\frac{\mathrm{d}u}{\mathrm{d}x}(x)\right)+q(x)u(x)-\lambda w(x)u(x)=f(x)
\end{equation}
以及下文将说明的适当边界条件. 作如下假设:\begin{enumerate}\item 函数 $p$, $p\prime$, $q$ 和 $w$ 在开区间 $(a,b)$ 上取实值且连续.\item 函数 $p$ 和 $w$ 在 $(a,b)$ 上为正.\end{enumerate}若 $a$ 和 $b$ 都有限, 且假设 1 和 2 在闭区间 $[a,b]$ 上成立, 则称 S-L 问题为\textbf{正则的}. 否则称该问题为\textbf{奇异的}.

形式地定义微分算子
\begin{equation}\label{eq:sturm-liouville-weighted-differential-operator}
Lu:=\dfrac{1}{w}[-(pu')'+qu],
\end{equation}
注意, 现在可以将 \eqref{eq:sturm-liouville-inhomogeneous-differential-equation} 写成
\begin{equation}\label{eq:sturm-liouville-weighted-operator-equation}
Lu-\lambda u=\dfrac{f}{w}.
\end{equation}
打算运用上面刚建立的理论, 将其作为从加权空间 $L_w^2(a,b)$ 到自身的算子的特征值问题加以分析. 但是, 奇异问题的分析侧重于不同于前述方法的其他方法, 因此只讨论正则问题.

这里使用加权空间 $L_w^2(a,b)$, 但对于正则问题, 这实际上不过是记号上的方便, 因为
\begin{equation}\label{eq:sturm-liouville-weighted-norm-equivalence}
\min_{x\in[a,b]}w(x)\|u\|_{L^2(a,b)}^2\leq\|u\|_{L_w^2(a,b)}^2\leq\max_{x\in[a,b]}w(x)\|u\|_{L^2(a,b)}^2,
\end{equation}
所以 $L^2$ 范数与 $L_w^2$ 范数等价. 将利用这一点为算子 $L$ 定义定义域.

考察应用中遇到的 S-L 问题最常见的一类边界条件, 即\textbf{非混合型}边界条件. 要求
\begin{equation}\label{eq:sturm-liouville-left-unmixed-boundary-condition}
\cos\alpha\,u(a)-\sin\alpha\,u'(a)=0,
\end{equation}

\begin{equation}\label{eq:sturm-liouville-right-unmixed-boundary-condition}
\cos\beta\,u(b)+\sin\beta\,u'(b)=0.
\end{equation}
现在定义定义域
\begin{equation}\label{eq:sturm-liouville-unmixed-operator-domain}
\mathcal{D}(L):=\{u\in H^2(a,b)\mid\eqref{eq:sturm-liouville-left-unmixed-boundary-condition}\text{ and }\eqref{eq:sturm-liouville-right-unmixed-boundary-condition}\text{ are satisfied}\}.
\end{equation}
下面证明以下定理.

\begin{Theorem}
\label{thm:regular-sturm-liouville-first-spectral-properties}
设 $(\mathcal{D}(L),L)$ 由 \eqref{eq:sturm-liouville-weighted-differential-operator} 和 \eqref{eq:sturm-liouville-unmixed-operator-domain} 定义. 以下结论成立:\begin{enumerate}\item $(\mathcal{D}(L),L)$ 的特征值为实数.\item $(\mathcal{D}(L),L)$ 的特征值以某个常数 $\lambda_G\in\mathbb{R}$ 为下界.\item 不同特征值所对应的特征函数在 $L_w^2(a,b)$ 中相互正交.\item 每个特征值的重数均为一.\end{enumerate}
\end{Theorem}

\begingroup
\postdisplaypenalty=10000
\renewcommand{\endProof}{\par\nobreak\hfill$\blacksquare$\par\vspace{0.1cm}}
\begin{Proof}
首先用分部积分证明 Lagrange 恒等式; 即对于 $H^2(a,b)$ 中的任意 $u$ 和 $v$, 有
\begin{equation}\label{eq:sturm-liouville-lagrange-identity}
(Lu,v)_w-(u,Lv)_w=\left.p(x)[u(x)v'(x)-u'(x)v(x)]\right|_a^b.
\end{equation}
因而, 若 $u$ 和 $v$ 属于 $\mathcal{D}(L)$, 则可利用边界条件 \eqref{eq:sturm-liouville-left-unmixed-boundary-condition} 和 \eqref{eq:sturm-liouville-right-unmixed-boundary-condition} 得到
\begin{equation}\label{eq:sturm-liouville-domain-symmetry}
(Lu,v)_w=(u,Lv)_w,
\end{equation}
这证明了 $(\mathcal{D}(L),L)$ 对称. 因此, Theorem~\ref{thm:symmetric-operator-spectral-properties} 立即给出第 1 和第 3 条.

为证明第 2 条, 对所有 $u\in\mathcal{D}(L)$ 证明如下形式的能量估计:
\begin{equation}\label{eq:sturm-liouville-lower-energy-bound}
(Lu,u)_w\geq\lambda_G\|u\|_{L_w^2(a,b)}^2
\end{equation}
(这是椭圆 PDE 中 Gårding 不等式的类似形式 (参见 Section~\ref{subsec:elliptic-garding-inequality}), 因而采用记号 $\lambda_G$.) 仅在 $\tan\alpha,\tan\beta\in[0,\infty)$ 的情形下证明, 其余情形的证明留给读者 (Problem~\ref{ex:sturm-liouville-general-energy-bound}). 对任意 $u\in\mathcal{D}(L)$, 有\[\begin{aligned}
(Lu,u)&=\int_a^b-(pu')'u+q|u|^2\,\mathrm{d}x\\
&=\int_a^bp|u'|^2+q|u|^2\,\mathrm{d}x+p(a)u'(a)u(a)-p(b)u'(b)u(b)\\
&=\int_a^bp|u'|^2+q|u|^2\,\mathrm{d}x\\
&\quad+p(a)\tan\alpha|u'(a)|^2+p(b)\tan\beta|u'(b)|^2\\
&\geq\int_a^bq(x)|u(x)|^2\,\mathrm{d}x\\
&\geq\dfrac{\min_{x\in[a,b]}q(x)}{\max_{x\in[a,b]}w(x)}\|u\|_{L_w^2(a,b)}^2.
\end{aligned}\]只需注意到, 对于任意特征值 $\lambda$, 有
\begin{equation}\label{eq:sturm-liouville-eigenfunction-energy-identity}
(Lu,u)_w=(\lambda wu,u)=\lambda(u,u)_w.
\end{equation}
所以 $\lambda\geq\lambda_G$, 这就得到第 2 条.

第 4 条由 ODE 初值问题的唯一性定理立即得到, 该定理表明, 边界条件 \eqref{eq:sturm-liouville-left-unmixed-boundary-condition} 或 \eqref{eq:sturm-liouville-right-unmixed-boundary-condition} 中的任意一个都在相差一个乘法常数的意义下确定齐次 ODE $Lu=\lambda u$ 的解.
\end{Proof}
\endgroup
